\subsection{Cardinality without AC}
In a world without AC, one must accept that cardinalities themselves are not totally ordered. 
This is because there is no way to actually prove that given two sets $A$ and $B$, one of the
two is injectable into the other via some injective function. Nonetheless, we can still work 
with the concept of "successor-cardinality". 

\begin{definition}[Successor cardinality]
    The successor cardinality $k^{+}$ of a given cardinality $k$ is 
    the smallest possible cardinality strictly greater than $k$ under cardinality-comparison. 
\end{definition}

\begin{lemma}[Existence of the successor cardinality]
    Every cardinality has a successor.
\end{lemma}
\begin{proof}
    The Cantor's theorem (theorem \ref{thm:Cantor}) states that for every set $A$, 
    the power set of $A$ is strictly greater than $A$ itself. Hence, no cardinality 
    can ever be the greatest one, since the cardinality of every set within it will 
    be strictly greater. \\
\end{proof}
\newpage